% ═══════════════════════════════════════════════════════════════
% SLIDES: Objetos de clase & hay
% Klasse 7 | Unidad 2 | Mi instituto
% ═══════════════════════════════════════════════════════════════

\documentclass[aspectratio=169, 11pt]{beamer}

\usepackage[utf8]{inputenc}
\usepackage[T1]{fontenc}
\usepackage[ngerman]{babel}
\usepackage{palatino}
\usepackage{helvet}
\usepackage{tabularx}
\usepackage{booktabs}
\usepackage{array}
\usepackage{tikz}
\usepackage{amssymb}

% ═══════════════════════════════════════════════════════════════
% THEME & FARBEN
% ═══════════════════════════════════════════════════════════════

\usetheme{default}
\usecolortheme{default}

% Spanisch-Farben
\definecolor{spanischrot}{HTML}{C41E3A}
\definecolor{dunkelgrau}{HTML}{333333}
\definecolor{hellgrau}{HTML}{F5F5F5}
\definecolor{erfolg}{HTML}{2E7D32}
\definecolor{fehler}{HTML}{C62828}

% Beamer-Anpassungen
\setbeamercolor{frametitle}{fg=white, bg=spanischrot}
\setbeamercolor{title}{fg=white, bg=spanischrot}
\setbeamercolor{structure}{fg=spanischrot}
\setbeamercolor{block title}{fg=white, bg=spanischrot}
\setbeamercolor{block body}{fg=dunkelgrau, bg=hellgrau}
\setbeamercolor{item}{fg=spanischrot}
\setbeamercolor{subitem}{fg=spanischrot}

\setbeamerfont{frametitle}{series=\bfseries, size=\Large}
\setbeamerfont{title}{series=\bfseries, size=\huge}

% Keine Navigation
\setbeamertemplate{navigation symbols}{}
\setbeamertemplate{footline}[frame number]

% ═══════════════════════════════════════════════════════════════
% BEFEHLE
% ═══════════════════════════════════════════════════════════════

% Spanisch hervorheben
\newcommand{\es}[1]{\textbf{\textcolor{spanischrot}{#1}}}
% Deutsch (klein, grau)
\newcommand{\de}[1]{{\small\textcolor{dunkelgrau}{\textit{#1}}}}
% Richtig/Falsch
\newcommand{\richtig}[1]{\textcolor{erfolg}{\checkmark\,#1}}
\newcommand{\falsch}[1]{\textcolor{fehler}{$\times$\,#1}}

% ═══════════════════════════════════════════════════════════════
% DOKUMENT
% ═══════════════════════════════════════════════════════════════

\begin{document}

% ---------------------------------------------------------------
% TITELFOLIE
% ---------------------------------------------------------------
{
\setbeamercolor{background canvas}{bg=spanischrot}
\begin{frame}[plain]
\begin{center}
\vspace{2cm}
{\Huge\bfseries\textcolor{white}{Objetos de clase}}\\[0.5cm]
{\Large\textcolor{white}{Was gibt es im Klassenraum?}}\\[1.5cm]
{\large\textcolor{white}{Spanisch · Klasse 7 · Unidad 2}}
\end{center}
\end{frame}
}

% ---------------------------------------------------------------
% LERNZIELE
% ---------------------------------------------------------------
\begin{frame}{Heute lernst du...}
\begin{itemize}
    \item[\textcolor{spanischrot}{\textbullet}] \textbf{Vokabeln:} Gegenstände im Klassenraum \& Schulsachen
    \item[\textcolor{spanischrot}{\textbullet}] \textbf{Grammatik:} \es{hay} = es gibt
    \item[\textcolor{spanischrot}{\textbullet}] \textbf{Kommunikation:} Beschreiben, was es in einem Raum gibt
\end{itemize}

\vspace{1cm}

\begin{block}{Am Ende kannst du...}
...deinen Klassenraum auf Spanisch beschreiben!
\end{block}
\end{frame}

% ---------------------------------------------------------------
% VOKABELN 1: IM KLASSENRAUM
% ---------------------------------------------------------------
\begin{frame}{En el aula \de{-- Im Klassenraum}}

\begin{columns}[T]
\begin{column}{0.48\textwidth}
\begin{tabular}{@{}ll@{}}
\es{la mesa} & \de{der Tisch} \\[0.5em]
\es{la silla} & \de{der Stuhl} \\[0.5em]
\es{la pizarra} & \de{die Tafel} \\[0.5em]
\es{la ventana} & \de{das Fenster} \\[0.5em]
\es{la puerta} & \de{die Tür} \\
\end{tabular}
\end{column}

\begin{column}{0.48\textwidth}
\begin{tabular}{@{}ll@{}}
\es{la pared} & \de{die Wand} \\[0.5em]
\es{la papelera} & \de{der Papierkorb} \\[0.5em]
\es{el reloj} & \de{die Uhr} \\[0.5em]
\es{el ordenador} & \de{der Computer} \\[0.5em]
\es{el mapa} & \de{die Landkarte} \\
\end{tabular}
\end{column}
\end{columns}

\vspace{0.8cm}

\begin{alertblock}{Achtung: Ausnahme!}
\es{el mapa} -- endet auf \textbf{-a}, ist aber \textbf{maskulin}!
\end{alertblock}

\end{frame}

% ---------------------------------------------------------------
% VOKABELN 2: SCHULSACHEN
% ---------------------------------------------------------------
\begin{frame}{En la mochila \de{-- Im Rucksack}}

\begin{columns}[T]
\begin{column}{0.48\textwidth}
\begin{tabular}{@{}ll@{}}
\es{la mochila} & \de{der Rucksack} \\[0.5em]
\es{el libro} & \de{das Buch} \\[0.5em]
\es{el cuaderno} & \de{das Heft} \\[0.5em]
\es{el bolígrafo} & \de{der Kugelschreiber} \\
\end{tabular}
\end{column}

\begin{column}{0.48\textwidth}
\begin{tabular}{@{}ll@{}}
\es{el lápiz} & \de{der Bleistift} \\[0.5em]
\es{el estuche} & \de{die Federtasche} \\[0.5em]
\es{la carpeta} & \de{die Mappe} \\[0.5em]
\es{la calculadora} & \de{der Taschenrechner} \\
\end{tabular}
\end{column}
\end{columns}

\vspace{1cm}

\begin{block}{Tipp}
\es{el bolígrafo} $\rightarrow$ kurz: \es{el boli} \de{(der Kuli)}
\end{block}

\end{frame}

% ---------------------------------------------------------------
% GRAMMATIK: HAY
% ---------------------------------------------------------------
\begin{frame}{\es{hay} = es gibt}

\begin{center}
{\Huge \es{hay}}\\[0.3cm]
{\Large = es gibt}
\end{center}

\vspace{0.5cm}

\begin{block}{Wichtig}
\es{hay} ist \textbf{unveränderlich} -- Singular und Plural gleich!
\end{block}

\vspace{0.5cm}

\begin{columns}[T]
\begin{column}{0.48\textwidth}
\textbf{Verwendung:}\\[0.3em]
\es{hay} + un/una\\
\es{hay} + Zahl\\
\es{hay} + muchos/muchas
\end{column}

\begin{column}{0.48\textwidth}
\textbf{Beispiele:}\\[0.3em]
\es{Hay una mesa.}\\
\es{Hay veinticinco sillas.}\\
\es{Hay muchos libros.}
\end{column}
\end{columns}

\end{frame}

% ---------------------------------------------------------------
% TYPISCHE FEHLER
% ---------------------------------------------------------------
\begin{frame}{Typische Fehler vermeiden!}

\begin{center}
{\Large Nach \es{hay} kommt \textbf{KEIN} bestimmter Artikel!}
\end{center}

\vspace{1cm}

\begin{columns}[c]
\begin{column}{0.45\textwidth}
\begin{center}
{\Large\textcolor{fehler}{FALSCH}}\\[0.5em]
\falsch{*Hay \textit{el} libro.}\\[0.3em]
\falsch{*Hay \textit{la} mesa.}\\[0.3em]
\falsch{*Es hay una mesa.}
\end{center}
\end{column}

\begin{column}{0.45\textwidth}
\begin{center}
{\Large\textcolor{erfolg}{RICHTIG}}\\[0.5em]
\richtig{Hay \textbf{un} libro.}\\[0.3em]
\richtig{Hay \textbf{una} mesa.}\\[0.3em]
\richtig{Hay una mesa.}
\end{center}
\end{column}
\end{columns}

\end{frame}

% ---------------------------------------------------------------
% FRAGEN MIT HAY
% ---------------------------------------------------------------
\begin{frame}{Preguntas con \es{hay}}

\begin{tabular}{@{}ll@{}}
\es{¿Qué hay en la clase?} & \de{Was gibt es in der Klasse?} \\[0.8em]
\es{¿Qué hay en tu mochila?} & \de{Was gibt es in deinem Rucksack?} \\[0.8em]
\es{¿Cuántos libros hay?} & \de{Wie viele Bücher gibt es?} \\[0.8em]
\es{¿Cuántas sillas hay?} & \de{Wie viele Stühle gibt es?} \\
\end{tabular}

\vspace{1cm}

\begin{block}{Merke}
\es{cuánt\textbf{o}s} (maskulin) vs. \es{cuánt\textbf{a}s} (feminin)\\[0.3em]
$\rightarrow$ passt sich dem Nomen an!
\end{block}

\end{frame}

% ---------------------------------------------------------------
% BEISPIELE
% ---------------------------------------------------------------
\begin{frame}{Ejemplos \de{-- Beispiele}}

\begin{itemize}
    \item[\textcolor{spanischrot}{\textbullet}] \es{En mi clase hay una pizarra grande.}
    \item[\textcolor{spanischrot}{\textbullet}] \es{Hay veinticinco mesas y veinticinco sillas.}
    \item[\textcolor{spanischrot}{\textbullet}] \es{En mi mochila hay dos libros y tres cuadernos.}
    \item[\textcolor{spanischrot}{\textbullet}] \es{En la pared hay un mapa de España.}
\end{itemize}

\vspace{1cm}

\begin{block}{Jetzt du!}
Beschreibe deinen Klassenraum mit \es{hay}!
\end{block}

\end{frame}

% ---------------------------------------------------------------
% ÜBUNG
% ---------------------------------------------------------------
\begin{frame}{Übung: $\rightarrow$ Arbeitsblatt}

\begin{center}
{\Large Bearbeite jetzt das \textbf{Arbeitsblatt}!}
\end{center}

\vspace{1cm}

\begin{block}{Aufgaben}
\begin{enumerate}
    \item Vokabeln zuordnen
    \item \es{el} oder \es{la} ergänzen
    \item Sätze mit \es{hay} bilden
    \item Lückentext ausfüllen
    \item Fragen beantworten
\end{enumerate}
\end{block}

\vspace{0.5cm}

\begin{center}
{\small\textcolor{dunkelgrau}{Bei Fragen: Hand heben!}}
\end{center}

\end{frame}

% ---------------------------------------------------------------
% ZUSAMMENFASSUNG
% ---------------------------------------------------------------
\begin{frame}{Zusammenfassung}

\begin{block}{Das hast du heute gelernt:}
\begin{itemize}
    \item Vokabeln: Klassenraum + Schulsachen
    \item \es{hay} = es gibt (unveränderlich!)
    \item Nach \es{hay} kein bestimmter Artikel (el/la)
    \item Fragen: \es{¿Qué hay...?} / \es{¿Cuántos/as...?}
\end{itemize}
\end{block}

\vspace{0.8cm}

\begin{center}
{\Large\textcolor{spanischrot}{¡Muy bien!}}
\end{center}

\end{frame}

% ---------------------------------------------------------------
% SCHLUSSFOLIE
% ---------------------------------------------------------------
{
\setbeamercolor{background canvas}{bg=spanischrot}
\begin{frame}[plain]
\begin{center}
\vspace{3cm}
{\Huge\bfseries\textcolor{white}{¡Muy bien!}}\\[1cm]
{\Large\textcolor{white}{Bis zur nächsten Stunde!}}
\end{center}
\end{frame}
}

\end{document}
